\section{Stochastic Dispatch and Market-Boundary Optimization}
\label{sec:stochastic-dispatch-summary}

\subsection{Problem Form and Runtime Scope}

This notebook chapter summarizes the stochastic market-boundary optimization
material maintained separately in the standalone document
\texttt{sections/17\_two\_stage\_stochastic.tex}.  The underlying formulation is
a multi-stage stochastic dispatch and maintenance-planning model with scenario
tree structure, expected-cost and risk terms, and decomposition-oriented solve
strategies.

\subsection{Code Map}

The material summarized here is aligned primarily with:
\begin{itemize}
  \item \texttt{src/analysis/market\_simulation.cpp} for market-clearing and scheduling-side model infrastructure,
  \item the MILP solver layer documented in Section~\ref{sec:engine-solver-system},
  \item the standalone research-style derivation in \texttt{sections/17\_two\_stage\_stochastic.tex} and its compiled PDF companion.
\end{itemize}

\subsection{Pseudocode Map}

The detailed technical notebook pseudocode already covers the main deterministic
building blocks reused here:
\begin{itemize}
  \item Algorithm A16 for DC OPF-type operational subproblems,
  \item Algorithm A20 for hierarchical annual and chronological dispatch orchestration,
  \item Algorithm E9 and Section~\ref{sec:native-branch-and-cut} for MILP solve mechanics.
\end{itemize}
The full stochastic extensive-form and decomposition derivations remain in the
standalone stochastic-dispatch manuscript rather than being duplicated inside
this notebook.

\subsection{Current Status}

The stochastic-dispatch material in this repository is best understood as a
hybrid of implementation-facing documentation and a standalone methods paper.
For the notebook build, the standalone paper is kept external and this section
acts as the normalized notebook-facing summary chapter.

\subsection{Canonical Problem Form}

In condensed form, the stochastic planning model can be read as
\begin{align}
\min_{x,\,y(\omega)} \quad & c^\top x + \sum_{\omega\in\Omega} p_\omega\, q_\omega^\top y(\omega) + \text{risk terms} \\
\text{s.t.}\quad & A x \le b, \\
& T_\omega x + W_\omega y(\omega) \le h_\omega, \qquad \forall \omega\in\Omega,
\end{align}
with first-stage boundary or maintenance decisions in \(x\) and scenario-wise
recourse decisions in \(y(\omega)\).

\subsection{How This Chapter Relates to the Standalone Manuscript}

The standalone file \texttt{sections/17\_two\_stage\_stochastic.tex} is still the
full derivation and should continue to be compiled on its own when the research
paper version is needed.  The technical notebook now keeps only the normalized
summary chapter so the main notebook build remains a single coherent document
instead of recursively including a second document preamble.